\label{mab:rm:ch:chaos}

\section{Construcci\'on intr\'inseca de la familia dodecaf\'asica}

Las constantes de Feigenbaum caracterizan la universalidad de la
duplicaci\'on de periodo. Su derivaci\'on en HMT exige construir previamente
una familia anal\'itica unimodal, demostrar su criticidad cuadr\'atica y
estudiar posteriormente la acci\'on del operador de renormalizaci\'on.

El mapa inicial no procede de una elecci\'on arm\'onica externa. El retorno
de \(\TPK\) determina la rueda dodecaf\'asica orientada
\(U_{12}^{\rm sign}\). En el sector uniforme \(\eta=0\), la traza
normalizada de sus doce sectores fija \(w_m=1/12\), y el funcional de fase
\begin{equation}
\mathcal P_{\rm crit}:U_{12}^{\rm sign}\longrightarrow\R/2\pi\Z,
\qquad
m\longmapsto \frac{(30m-10)\pi}{180}
=\pi\frac{3m-1}{18}
\label{mab:rm:eq:critical-phase-reader}
\end{equation}
conserva orden y orientaci\'on. La composici\'on causal es
\[
\APP\to\TRIT\to\TPK\to U_{12}^{\rm sign}
\xrightarrow{\mathcal P_{\rm crit}}
\{(\phi_m,w_m)\}_{m=1}^{12}
\to u_0\to f_\lambda.
\]
Toda modificaci\'on de \(\mathcal P_{\rm crit}\), de los pesos o del orden de
las fases altera la familia y constituye un criterio de refutaci\'on. Estas
elecciones deben fijarse con anterioridad al contraste con \(\delta_F\).

La familia dodecaf\'asica resultante parte de

\[
a_m=\frac{3m-1}{18},\qquad
\phi_m=\pi a_m,\qquad m=1,\ldots,12,
\]

y del segundo momento

\[
D_0=\frac1{12}\sum_{m=1}^{12}a_m^2.
\]

La identidad

\begin{equation}
\sum_{m=1}^{12}(3m-1)^2=5394=6\cdot899
\label{mab:rm:eq:chaos-sum2}
\end{equation}

produce

\[
D_0=\frac{899}{648},\qquad
s_0=\frac{36}{\pi\sqrt{899}}.
\]

Por tanto,

\begin{equation}
s_0\phi_m=\frac{2(3m-1)}{\sqrt{899}}.
\label{mab:rm:eq:pi-cancellation}
\end{equation}

El factor \(\pi\) se cancela antes de efectuar las iteraciones. La familia no
depende de \(\delta_F\) ni de \(\alpha_F\) como datos de entrada.

La cancelaci\'on es una propiedad estructural de la normalizaci\'on. Antes de
normalizar,
el segundo momento angular es
\[
\frac1{12}\sum_{m=1}^{12}\phi_m^2
=\pi^2D_0.
\]
La condici\'on de criticidad cuadr\'atica fija
\(s_0^2\pi^2D_0=2\), y por ello cada fase normalizada depende s\'olo de la
aritm\'etica entera de \(3m-1\). El funcional conserva la procedencia angular de
la rueda, pero la iteraci\'on recibe una familia real intr\'inseca. Esto es
precisamente lo que permite distinguir una derivaci\'on de una calibraci\'on
por \(\pi\).

\section{Representaci\'on anal\'itica exacta y expresi\'on trigonom\'etrica cerrada}

Definimos

\begin{equation}
u_0(x)=1-\frac1{12}\sum_{m=1}^{12}
\cos\!\left(\frac{2(3m-1)}{\sqrt{899}}x\right).
\label{mab:rm:eq:chaos-profile}
\end{equation}

El desarrollo en el origen es

\begin{equation}
u_0(x)=x^2-\frac{715769}{2424603}x^4+O(x^6).
\label{mab:rm:eq:chaos-taylor}
\end{equation}

La suma finita admite la siguiente expresi\'on cerrada. Si \(y=x/\sqrt{899}\),

\begin{equation}
u_0(x)=1-\frac{\cos(37y)\sin(36y)}{12\sin(3y)}
=1-\frac{\sin(73y)-\sin y}{24\sin(3y)}.
\label{mab:rm:eq:chaos-closed}
\end{equation}

La singularidad aparente en \(y=0\) es removible y reproduce \cref{mab:rm:eq:chaos-taylor}.

\begin{proof}[Derivaci\'on de la forma cerrada]
Los doce argumentos forman una progresi\'on aritm\'etica:
\[
2(3m-1)y=(6m-2)y,\qquad m=1,\ldots,12.
\]
La suma de cosenos de una progresi\'on satisface
\[
\sum_{m=1}^{N}\cos(a+md)
=\frac{\sin(Nd/2)}{\sin(d/2)}
 \cos\!\left(a+\frac{N+1}{2}d\right).
\]
Tomando \(N=12\), \(d=6y\) y \(a=-2y\) resulta
\[
\sum_{m=1}^{12}\cos((6m-2)y)
=\frac{\sin36y}{\sin3y}\cos37y.
\]
La identidad producto--a--suma
\(2\cos37y\sin36y=\sin73y-\sin y\) da la segunda expresi\'on de
\cref{mab:rm:eq:chaos-closed}.
\end{proof}

El coeficiente cu\'artico tampoco se ajusta. Las sumas enteras
\begin{equation}
\sum_{m=1}^{12}(3m-1)^2=5394,\qquad
\sum_{m=1}^{12}(3m-1)^4=4294614
\label{mab:rm:eq:chaos-moments24}
\end{equation}
insertadas en el desarrollo de \(\cos z\) producen exactamente
\[
\frac1{24\cdot12}\sum_m
\left(\frac{2(3m-1)}{\sqrt{899}}\right)^4
=\frac{715769}{2424603}.
\]
As\'i, tanto la normalizaci\'on cuadr\'atica como la primera correcci\'on no
lineal son momentos finitos de la misma rueda.

\section{Pertenencia a la clase anal\'itica unimodal cuadr\'atica}

Consideremos

\[
f_\lambda(x)=1-\lambda u_0(x),\qquad \lambda>0.
\]

Para \(0<x\le1\), todos los argumentos

\[
\frac{2(3m-1)}{\sqrt{899}}x
\]

pertenecen a \((0,\pi)\). Por ello,

\[
u_0'(x)>0,\qquad u_0'''(x)<0.
\]

Como \(u_0\) es par, \(x=0\) es su único punto cr\'itico y \(u_0''(0)=2\). El Schwarziano de \(f_\lambda\) es

\begin{equation}
S(f_\lambda)
=\frac{u_0'''}{u_0'}
-\frac32\left(\frac{u_0''}{u_0'}\right)^2<0
\qquad(0<|x|\le1).
\label{mab:rm:eq:negative-schwarzian}
\end{equation}

\begin{theorem}[Pertenencia exacta a la clase unimodal]
Para
\[
0<\lambda\le \frac{2}{u_0(1)},
\]
la familia \(f_\lambda=1-\lambda u_0\) es una autoaplicaci\'on real--anal\'itica de \([-1,1]\), par, con un único punto cr\'itico cuadr\'atico en el origen y Schwarziano negativo fuera de ese punto. Por tanto pertenece, sin usar valores de Feigenbaum como valor objetivo, a la clase \(S\)-unimodal cuadr\'atica sobre la que act\'ua el renormalizador de duplicaci\'on de periodo \cite{rm:Feigenbaum1978,rm:Lanford1982}. Para \(\lambda\) fuera de ese intervalo se conserva el germen anal\'itico y su forma local, pero no se afirma autom\'aticamente la autoaplicaci\'on del intervalo.
\end{theorem}

El teorema establece la contribuci\'on intr\'inseca de HMT previa al teorema
de universalidad: el germen anal\'itico y su tipo de criticidad.

\begin{proof}[Detalles del signo y de la autoaplicaci\'on]
Derivando \cref{mab:rm:eq:chaos-profile},
\begin{align*}
u_0'(x)&=\frac1{12}\sum_m k_m\sin(k_mx),\\
u_0'''(x)&=-\frac1{12}\sum_m k_m^3\sin(k_mx),
\qquad k_m=\frac{2(3m-1)}{\sqrt{899}}.
\end{align*}
Para \(0<x\le1\), \(0<k_mx<\pi\), luego todos los sumandos de la primera
l\'inea son positivos y los de la segunda negativos. La paridad extiende la
monoton\'ia a ambos lados del origen. Como \(u_0(0)=0\) y
\(0\le u_0(x)\le u_0(1)\), la cota
\(0<\lambda\le2/u_0(1)\) implica
\(-1\le1-\lambda u_0(x)\le1\). Finalmente,
\(u_0''(0)=1/12\sum k_m^2=2\), de modo que el punto cr\'itico es
cuadr\'atico.
\end{proof}

\section{Sucesiones finitas de par\'ametros superestables}

La sucesi\'on superestable certificada hasta nivel ocho produce

\[
\delta_8=4.66921218915\ldots,
\qquad
\alpha_8=2.50296847988\ldots.
\]

Su estatuto es \emph{Certificado finito libre de valores objetivo}: cada
aproximante se obtiene de la familia previamente definida, sin introducir
como dato el l\'imite universal. La coincidencia decimal proporciona un
indicador de convergencia, pero no sustituye el teorema de hiperbolicidad.

Una extensi\'on num\'erica independiente de la misma sucesi\'on, con paso
adaptado a la separaci\'on entre ra\'ices y control del periodo exacto,
alcanza provisionalmente:

\begin{center}
\small
\begin{tabular}{@{}ccll@{}}
\toprule
Nivel & \(\lambda_n\) & \(\delta_n\) & \(\alpha_n\)\\
\midrule
11 & 1.87403828195439908045 & 4.66920170694424982745 & 2.50290847517165181764\\
12 & 1.87403836411023460451 & 4.66920163036308158848 & 2.50290800379001546915\\
13 & 1.87403838170549870372 & 4.66920161361839219913 & 2.50290790268748193108\\
14 & 1.87403838547386545431 & 4.66920161007472591023 & 2.50290788100969754274\\
\bottomrule
\end{tabular}
\end{center}

Estos niveles poseen el estatuto de \emph{Certificado num\'erico nuevo
pendiente de registro reproducible}; no sustituyen el certificado persistente
hasta nivel ocho. Una extrapolaci\'on de Aitken proporciona
\(\lambda_\infty^{\HMT}\simeq1.8740383865008918080\) y
\(\alpha_F\simeq2.50290787509\) como candidatos de l\'imite, condicionados a
la demostraci\'on de hiperbolicidad y transversalidad.

Para evitar un procedimiento num\'erico no explicitado, el m\'etodo se formula como
sigue. Defina
\begin{equation}
F_n(\lambda)=f_\lambda^{\,2^n}(0).
\label{mab:rm:eq:superstable-equation}
\end{equation}
La ra\'iz superestable \(\lambda_n\) es la primera ra\'iz de \(F_n\) a la
derecha de \(\lambda_{n-1}\) que satisface simult\'aneamente
\[
f_{\lambda_n}^{\,2^j}(0)\ne0
\quad(0\le j<n).
\]
Esta desigualdad excluye los periodos divisores. Con intervalos disjuntos
\(I_n\ni\lambda_n\), se calculan
\begin{equation}
\delta_n=
\frac{\lambda_{n-1}-\lambda_{n-2}}
     {\lambda_n-\lambda_{n-1}},
\qquad
\alpha_n=-\frac{d_{n-1}}{d_n},
\label{mab:rm:eq:feigenbaum-approximants}
\end{equation}
donde \(d_n=f_{\lambda_n}^{\,2^{n-1}}(0)\) es la distancia orientada del
\'ultimo punto nuevo al punto cr\'itico. Un certificado persistente debe guardar
los intervalos de ra\'iz, la exclusi\'on de periodos menores y la propagaci\'on
de error de \cref{mab:rm:eq:feigenbaum-approximants}; un simple decimal impreso no
cumple esas tres condiciones.

La regla de discretizaci\'on constituye asimismo un control negativo. Una
cota inferior fija mayor que la separaci\'on
\(|\lambda_n-\lambda_{n-1}|\) puede omitir la primera ra\'iz y producir, no
obstante, una sucesi\'on num\'ericamente plausible. Por ello el incremento
debe escalar con la separaci\'on previa y cada candidato debe verificarse
mediante su periodo exacto.

\section{Condiciones de certificaci\'on del operador de renormalizaci\'on}

Sobre un espacio de funciones pares anal\'iticas, definimos

\begin{equation}
\mathcal R(f)(x)=-a(f)^{-1}f\bigl(f(-a(f)x)\bigr),
\qquad a(f)=-f(1).
\label{mab:rm:eq:renormalization}
\end{equation}

El cierre completo de \(\delta_F\) y \(\alpha_F\) requiere certificar conjuntamente:

\begin{enumerate}
\item existencia de un punto fijo \(g=\mathcal R(g)\) en una bola anal\'itica declarada;
\item unicidad local de ese punto fijo;
\item un único autovalor real inestable de \(D\mathcal R_g\), igual a \(\delta_F\);
\item radio espectral estrictamente menor que uno en el complemento estable;
\item transversalidad de la curva \(\lambda\mapsto f_\lambda\) respecto de la variedad estable de \(g\).
\end{enumerate}

Una implementaci\'on HMT natural emplea truncaciones de Chebyshev de orden
\(9^m\), aritm\'etica de intervalos y polinomios de radios. La prueba debe
proporcionar una bola, una inversa aproximada, una cota de defecto y una cota
de Lipschitz que satisfagan \(Y+Z(r)<r\). De este modo, la profundidad
non\'adica parametriza un certificado de punto fijo y no una mera sucesi\'on
de aproximaciones decimales.

En una base par de Chebyshev,
\[
g(x)=\sum_{j=0}^{N}g_{2j}T_{2j}(x),\qquad N=9^m,
\]
el residuo finito es \(F_N(g)=\Pi_N(\mathcal Rg-g)\). Si \(A_N\) aproxima
\((DF_N(\bar g))^{-1}\), las cantidades que deben certificarse son
\begin{align*}
Y&=\|A_NF_N(\bar g)\|,\\
Z_0&=\|I-A_NDF_N(\bar g)\|,\\
Z_1(r)&=\sup_{\|h\|\le r}
\|A_N(DF_N(\bar g+h)-DF_N(\bar g))\|.
\end{align*}
Una desigualdad
\(Y+(Z_0+Z_1(r))r<r\) encierra un punto fijo \`unico. El bloque espectral
posterior debe aislar una direcci\'on inestable y contraer el complemento. El
uso de \(9^m\) entra as\'i como ley de refinamiento del certificado, no como
decoraci\'on numerol\'ogica.

\begin{corollary}[Confluencia exacta del modo treinta]
El modo \(30\) admite tres realizaciones simult\'aneas y tipol\'ogicamente
distintas:
\begin{align*}
30&\longmapsto(3\cdot30,4\cdot30)=(90,120)
&&\text{en la completaci\'on constitutiva},\\
30&\longmapsto30^2-1=899=29\cdot31
&&\text{en la normalizaci\'on cr\'itica},\\
30&\longmapsto
M_{30}=\begin{pmatrix}30&899\\1&30\end{pmatrix}
&&\text{en la transformaci\'on hiperb\'olica de Pell}.
\end{align*}
Las tres identidades son exactas y se obtienen sin prescribir valores objetivo. El invariante conservado es
el entero central \(30\) junto con su orientaci\'on: extensi\'on \(3\to4\),
defecto cuadr\'atico \(-1\) y unidad de norma uno. Lo que no se afirma es la
igualdad entre el operador de vac\'io y el renormalizador de duplicaci\'on.
\end{corollary}

\begin{remark}[Control negativo]
Un segundo autovalor fuera del disco unidad, una bola de radios sin soluci\'on, p\'erdida de unicidad o tangencia de la familia HMT a la variedad estable refutan la promoci\'on universal. Ninguno refuta el germen exacto ni el teorema de Schwarziano.
\end{remark}

\begin{proposition}[Resultado principal]
La significaci\'on estructural de la constante de duplicaci\'on reside en la
cancelaci\'on dodecaf\'asica de \(\pi\), en la identidad
\(899=30^2-1\), en la forma trigonom\'etrica cerrada y en la pertenencia
demostrada a la clase anal\'itica sobre la que act\'ua el renormalizador. Su
expansi\'on decimal constituye exclusivamente la evaluaci\'on terminal.
\end{proposition}

\subsection*{Alcance lógico y dominio de validez}
Perfil, Taylor, normalizador, unidad de Pell, unimodalidad y Schwarziano: \emph{Teorema interno exacto}. Aproximantes hasta nivel ocho: \emph{Certificado finito persistente}. Niveles 11--14: \emph{Certificado num\'erico nuevo pendiente de durabilidad}. Hiperbolicidad y transversalidad del renormalizador: \emph{Programa abierto con criterio de refutación ejecutable}.

